\chapter{The Ext Functor}
\label{ch:v27}

The Ext functor is the derived correction to the left exact Hom functor. It classifies extensions in degree one and measures higher obstructions to lifting maps.

\section*{Learning goals}
After this chapter, the reader should be able to:
\begin{itemize}
\item compute Ext groups from projective resolutions and Hom complexes;
\item verify that \(\operatorname{Ext}^0\) recovers ordinary Hom;
\item use long exact Ext sequences and connecting morphisms in concrete examples;
\item interpret \(\operatorname{Ext}^1\) in terms of extensions where developed;
\item compute Ext for cyclic modules or quotients from short free resolutions;
\item distinguish variance in the first and second arguments and track induced maps correctly;
\end{itemize}
\section*{Conceptual roadmap}
\volroadmap
  {presentation or universal property}
  {exactness/localization}
  {structural consequence}

\section{Definition from projectives}
Choose a projective resolution $P_\bullet\to M$ and set $\operatorname{Ext}_A^n(M,N)=H^n(\operatorname{Hom}_A(P_\bullet,N))$.

\section{Degree zero}
$\operatorname{Ext}_A^0(M,N)\cong\operatorname{Hom}_A(M,N)$.

% BEGIN VOL05-EXPANSION V27-example-01
\begin{example}[Ext of cyclic groups]\label{ex:v27-expand-01}
Applying $\operatorname{Hom}_{\mathbb Z}(-,\mathbb Z/n)$ to the standard resolution of $\mathbb Z/m$ gives a cokernel of multiplication by $m$. Hence $\operatorname{Ext}^1_{\mathbb Z}(\mathbb Z/m,\mathbb Z/n)\cong\mathbb Z/\gcd(m,n)$.
\end{example}
% END VOL05-EXPANSION V27-example-01
\section{Long exact sequence}
Short exact sequences yield long exact Ext sequences.

\section{Projective vanishing}
Higher Ext out of a projective module is zero.

\section{Injective viewpoint}
Ext can also be computed using injective resolutions in the second variable.

% BEGIN VOL05-EXPANSION V27-example-02
\begin{example}[Projectives kill Ext]\label{ex:v27-expand-02}
If $P$ is projective, it has a length-zero projective resolution. Therefore $\operatorname{Ext}^i_R(P,N)=0$ for every $i>0$.
\end{example}
% END VOL05-EXPANSION V27-example-02
\section{Extension classes}
$\operatorname{Ext}^1(M,N)$ classifies short exact sequences $0\to N\to E\to M\to0$ up to the appropriate equivalence.

\section{Dimension shifting}
Syzygies reduce higher Ext to lower Ext.

% BEGIN VOL05-EXPANSION V27-example-03
\begin{example}[Extensions classified by Ext]\label{ex:v27-expand-03}
Elements of $\operatorname{Ext}^1_R(M,N)$ correspond to equivalence classes of short exact sequences $0\to N\to E\to M\to0$. The zero class is represented by the split extension.
\end{example}
% END VOL05-EXPANSION V27-example-03
\section{Local and graded variants}
Localization and graded structures refine Ext in common Noetherian settings.

\section{Duality preview}
Ext groups appear in duality, canonical modules, and deformation theory.

\section{Core structural results}

\begin{theorem}[Well-definedness of Ext]\label{thm:v27-01}
$H^n(\operatorname{Hom}(P_\bullet,N))$ is independent of the chosen projective resolution of $M$.
\begin{proof}
Comparison maps between resolutions are unique up to chain homotopy. Applying Hom turns chain homotopies into cochain homotopies, which induce identical cohomology maps.
\end{proof}
\end{theorem}

\begin{theorem}[Projective vanishing]\label{thm:v27-02}
If $P$ is projective, then $\operatorname{Ext}_A^n(P,N)=0$ for $n>0$.
\begin{proof}
Use the projective resolution concentrated in degree zero. The resulting Hom cochain complex has no positive cohomology.
\end{proof}
\end{theorem}

\begin{theorem}[Ext one classifies extensions]\label{thm:v27-03}
Equivalence classes of short exact sequences $0\to N\to E\to M\to0$ form a group naturally isomorphic to $\operatorname{Ext}_A^1(M,N)$.
\begin{proof}
Pull an extension back along a projective presentation of $M$ and choose a splitting over the projective term. The resulting map from the first syzygy to $N$ defines a degree-one cohomology class; changing the splitting changes it by a coboundary, and the construction is reversible via pushout.
\end{proof}
\end{theorem}

\section{Worked examples}

\begin{example}[Definition from projectives diagnostic]\label{ex:v27-worked-01}
Audit Definition from projectives in four stages: hypotheses, decisive mechanism, argument, and limitation. For the decisive mechanism, choose a projective/free resolution and compute cohomology after applying Hom. Choose a projective resolution $P_\bullet\to M$ and set $\operatorname{Ext}_A^n(M,N)=H^n(\operatorname{Hom}_A(P_\bullet,N))$. The decisive algebraic mechanism is to choose a projective/free resolution and compute cohomology after applying Hom. The stated hypotheses determine which structural conclusion is valid.
\end{example}

\begin{example}[Degree zero diagnostic]\label{ex:v27-worked-02}
Treat Degree zero as a short proof dossier. State the hypotheses explicitly, then choose a projective/free resolution and compute cohomology after applying Hom. Derive the conclusion and identify a failure mode if a required hypothesis is absent. $\operatorname{Ext}_A^0(M,N)\cong\operatorname{Hom}_A(M,N)$. The decisive algebraic mechanism is to choose a projective/free resolution and compute cohomology after applying Hom. The stated hypotheses determine which structural conclusion is valid.
\end{example}

\begin{example}[Long exact sequence diagnostic]\label{ex:v27-worked-03}
Audit Long exact sequence in four stages: hypotheses, decisive mechanism, argument, and limitation. For the decisive mechanism, choose a projective/free resolution and compute cohomology after applying Hom. Short exact sequences yield long exact Ext sequences. The decisive algebraic mechanism is to choose a projective/free resolution and compute cohomology after applying Hom. The stated hypotheses determine which structural conclusion is valid.
\end{example}

\begin{example}[Projective vanishing diagnostic]\label{ex:v27-worked-04}
For Projective vanishing, begin by isolating the assumptions. Next, choose a projective/free resolution and compute cohomology after applying Hom. Conclude by stating precisely what has been proved and where the reasoning can break down. Higher Ext out of a projective module is zero. The decisive algebraic mechanism is to choose a projective/free resolution and compute cohomology after applying Hom. The stated hypotheses determine which structural conclusion is valid.
\end{example}

\section{Solved problems}
Each extended problem combines several results from the chapter in a longer argument. It is intended to make hypothesis selection, proof strategy, and verification explicit before the shorter exercises begin.

\begin{problem}[Definition from projectives diagnostic]\label{prob:v27-problem-01}
Audit Definition from projectives in four stages: hypotheses, decisive mechanism, argument, and limitation. For the decisive mechanism, choose a projective/free resolution and compute cohomology after applying Hom.
\end{problem}
\begin{solution}
Choose a projective resolution $P_\bullet\to M$ and set $\operatorname{Ext}_A^n(M,N)=H^n(\operatorname{Hom}_A(P_\bullet,N))$. The decisive algebraic mechanism is to choose a projective/free resolution and compute cohomology after applying Hom. The stated hypotheses determine which structural conclusion is valid.
\end{solution}

\begin{problem}[Degree zero diagnostic]\label{prob:v27-problem-02}
Treat Degree zero as a short proof dossier. State the hypotheses explicitly, then choose a projective/free resolution and compute cohomology after applying Hom. Derive the conclusion and identify a failure mode if a required hypothesis is absent.
\end{problem}
\begin{solution}
$\operatorname{Ext}_A^0(M,N)\cong\operatorname{Hom}_A(M,N)$. The decisive algebraic mechanism is to choose a projective/free resolution and compute cohomology after applying Hom. The stated hypotheses determine which structural conclusion is valid.
\end{solution}

\begin{problem}[Long exact sequence diagnostic]\label{prob:v27-problem-03}
Audit Long exact sequence in four stages: hypotheses, decisive mechanism, argument, and limitation. For the decisive mechanism, choose a projective/free resolution and compute cohomology after applying Hom.
\end{problem}
\begin{solution}
Short exact sequences yield long exact Ext sequences. The decisive algebraic mechanism is to choose a projective/free resolution and compute cohomology after applying Hom. The stated hypotheses determine which structural conclusion is valid.
\end{solution}

\begin{problem}[Projective vanishing diagnostic]\label{prob:v27-problem-04}
For Projective vanishing, begin by isolating the assumptions. Next, choose a projective/free resolution and compute cohomology after applying Hom. Conclude by stating precisely what has been proved and where the reasoning can break down.
\end{problem}
\begin{solution}
Higher Ext out of a projective module is zero. The decisive algebraic mechanism is to choose a projective/free resolution and compute cohomology after applying Hom. The stated hypotheses determine which structural conclusion is valid.
\end{solution}

\begin{problem}[Injective viewpoint diagnostic]\label{prob:v27-problem-05}
Audit Injective viewpoint in four stages: hypotheses, decisive mechanism, argument, and limitation. For the decisive mechanism, choose a projective/free resolution and compute cohomology after applying Hom.
\end{problem}
\begin{solution}
Ext can also be computed using injective resolutions in the second variable. The decisive algebraic mechanism is to choose a projective/free resolution and compute cohomology after applying Hom. The stated hypotheses determine which structural conclusion is valid.
\end{solution}

\begin{problem}[Extension classes diagnostic]\label{prob:v27-problem-06}
Resolve Extension classes by separating assumptions from inference. State the assumptions, choose a projective/free resolution and compute cohomology after applying Hom, give the conclusion, and note one condition under which that conclusion is no longer justified.
\end{problem}
\begin{solution}
$\operatorname{Ext}^1(M,N)$ classifies short exact sequences $0\to N\to E\to M\to0$ up to the appropriate equivalence. The decisive algebraic mechanism is to choose a projective/free resolution and compute cohomology after applying Hom. The stated hypotheses determine which structural conclusion is valid.
\end{solution}

\begin{problem}[Dimension shifting diagnostic]\label{prob:v27-problem-07}
Resolve Dimension shifting by separating assumptions from inference. State the assumptions, choose a projective/free resolution and compute cohomology after applying Hom, give the conclusion, and note one condition under which that conclusion is no longer justified.
\end{problem}
\begin{solution}
Syzygies reduce higher Ext to lower Ext. The decisive algebraic mechanism is to choose a projective/free resolution and compute cohomology after applying Hom. The stated hypotheses determine which structural conclusion is valid.
\end{solution}

\begin{problem}[Local and graded variants diagnostic]\label{prob:v27-problem-08}
Treat Local and graded variants as a short proof dossier. State the hypotheses explicitly, then choose a projective/free resolution and compute cohomology after applying Hom. Derive the conclusion and identify a failure mode if a required hypothesis is absent.
\end{problem}
\begin{solution}
Localization and graded structures refine Ext in common Noetherian settings. The decisive algebraic mechanism is to choose a projective/free resolution and compute cohomology after applying Hom. The stated hypotheses determine which structural conclusion is valid.
\end{solution}

\begin{problem}[Well-definedness of Ext]\label{prob:v27-problem-09}
Prove the following structural statement and identify the hypothesis that makes the conclusion work: $H^n(\operatorname{Hom}(P_\bullet,N))$ is independent of the chosen projective resolution of $M$.
\end{problem}
\begin{solution}
Comparison maps between resolutions are unique up to chain homotopy. Applying Hom turns chain homotopies into cochain homotopies, which induce identical cohomology maps. The proof hinges on the following algebraic step: choose a projective/free resolution and compute cohomology after applying Hom.
\end{solution}

\begin{problem}[Projective vanishing]\label{prob:v27-problem-10}
Prove the following structural statement and identify the hypothesis that makes the conclusion work: If $P$ is projective, then $\operatorname{Ext}_A^n(P,N)=0$ for $n>0$.
\end{problem}
\begin{solution}
Use the projective resolution concentrated in degree zero. The resulting Hom cochain complex has no positive cohomology. The proof hinges on the following algebraic step: choose a projective/free resolution and compute cohomology after applying Hom.
\end{solution}

\begin{problem}[Ext one classifies extensions]\label{prob:v27-problem-11}
Prove the following structural statement and identify the hypothesis that makes the conclusion work: Equivalence classes of short exact sequences $0\to N\to E\to M\to0$ form a group naturally isomorphic to $\operatorname{Ext}_A^1(M,N)$.
\end{problem}
\begin{solution}
Pull an extension back along a projective presentation of $M$ and choose a splitting over the projective term. The resulting map from the first syzygy to $N$ defines a degree-one cohomology class; changing the splitting changes it by a coboundary, and the construction is reversible via pushout. The proof hinges on the following algebraic step: choose a projective/free resolution and compute cohomology after applying Hom.
\end{solution}

\begin{problem}[Chapter synthesis]\label{prob:v27-problem-12}
Connect the definitions and principal structural theorems of this chapter into one reusable commutative-algebra strategy.
\end{problem}
\begin{solution}
The Ext functor is the derived correction to the left exact Hom functor. It classifies extensions in degree one and measures higher obstructions to lifting maps. A chapter-specific strategy is to choose a projective/free resolution and compute cohomology after applying Hom; consequently, Ext measures the failure of Hom to be exact and classifies extension phenomena in low degree.
\end{solution}

\section{Exercises with complete solutions}

\begin{exercise}\label{exr:v27-01}
For Definition from projectives, distinguish the formal setup from the decisive step. First record the hypotheses; then choose a projective/free resolution and compute cohomology after applying Hom. Explain the conclusion and a failure mode outside those hypotheses.
\end{exercise}
\begin{hint}
Choose a projective resolution \(P_\bullet\to M\) and apply \(\operatorname{Hom}_A(P_\bullet,N)\); Ext is the cohomology of that cochain complex.
\end{hint}
\begin{solution}
Choose a projective resolution $P_\bullet\to M$ and set $\operatorname{Ext}_A^n(M,N)=H^n(\operatorname{Hom}_A(P_\bullet,N))$. In this chapter the concrete consequence is that Ext measures the failure of Hom to be exact and classifies extension phenomena in low degree.
\end{solution}

\begin{exercise}\label{exr:v27-02}
Work through Degree zero from the assumptions outward. Under the stated hypotheses, choose a projective/free resolution and compute cohomology after applying Hom. Give the resulting conclusion and one boundary case where the argument no longer applies.
\end{exercise}
\begin{hint}
Degree zero cohomology recovers \(\operatorname{Hom}_A(M,N)\); verify this from the first two terms of the resolution.
\end{hint}
\begin{solution}
$\operatorname{Ext}_A^0(M,N)\cong\operatorname{Hom}_A(M,N)$. In this chapter the concrete consequence is that Ext measures the failure of Hom to be exact and classifies extension phenomena in low degree.
\end{solution}

\begin{exercise}\label{exr:v27-03}
Determine which hypotheses are actually needed for Long exact sequence. Then choose a projective/free resolution and compute cohomology after applying Hom. State the conclusion and identify what can fail when one essential hypothesis is removed.
\end{exercise}
\begin{hint}
For a cyclic quotient, apply Hom to a short free resolution and compute kernel/cokernel of multiplication.
\end{hint}
\begin{solution}
Short exact sequences yield long exact Ext sequences. In this chapter the concrete consequence is that Ext measures the failure of Hom to be exact and classifies extension phenomena in low degree.
\end{solution}

\begin{exercise}\label{exr:v27-04}
Determine which hypotheses are actually needed for Projective vanishing. Then choose a projective/free resolution and compute cohomology after applying Hom. State the conclusion and identify what can fail when one essential hypothesis is removed.
\end{exercise}
\begin{hint}
A short exact sequence yields a long exact Ext sequence; mark variance so arrows go in the correct direction.
\end{hint}
\begin{solution}
Higher Ext out of a projective module is zero. In this chapter the concrete consequence is that Ext measures the failure of Hom to be exact and classifies extension phenomena in low degree.
\end{solution}

\begin{exercise}\label{exr:v27-05}
For Injective viewpoint, distinguish the formal setup from the decisive step. First record the hypotheses; then choose a projective/free resolution and compute cohomology after applying Hom. Explain the conclusion and a failure mode outside those hypotheses.
\end{exercise}
\begin{hint}
\(\operatorname{Ext}^1\) classifies extensions in the chapter's setting; identify the split extension as the zero class.
\end{hint}
\begin{solution}
Ext can also be computed using injective resolutions in the second variable. In this chapter the concrete consequence is that Ext measures the failure of Hom to be exact and classifies extension phenomena in low degree.
\end{solution}

\begin{exercise}\label{exr:v27-06}
Determine which hypotheses are actually needed for Extension classes. Then choose a projective/free resolution and compute cohomology after applying Hom. State the conclusion and identify what can fail when one essential hypothesis is removed.
\end{exercise}
\begin{hint}
Projective first arguments have vanishing higher Ext when computed projectively.
\end{hint}
\begin{solution}
$\operatorname{Ext}^1(M,N)$ classifies short exact sequences $0\to N\to E\to M\to0$ up to the appropriate equivalence. In this chapter the concrete consequence is that Ext measures the failure of Hom to be exact and classifies extension phenomena in low degree.
\end{solution}

\begin{exercise}\label{exr:v27-07}
For Dimension shifting, distinguish the formal setup from the decisive step. First record the hypotheses; then choose a projective/free resolution and compute cohomology after applying Hom. Explain the conclusion and a failure mode outside those hypotheses.
\end{exercise}
\begin{hint}
Changing the second argument is covariant while changing the first is contravariant; write induced maps before composing them.
\end{hint}
\begin{solution}
Syzygies reduce higher Ext to lower Ext. In this chapter the concrete consequence is that Ext measures the failure of Hom to be exact and classifies extension phenomena in low degree.
\end{solution}

\begin{exercise}\label{exr:v27-08}
For Local and graded variants, distinguish the formal setup from the decisive step. First record the hypotheses; then choose a projective/free resolution and compute cohomology after applying Hom. Explain the conclusion and a failure mode outside those hypotheses.
\end{exercise}
\begin{hint}
Independence of the chosen resolution is a homotopy-invariance statement, not equality of the resulting cochain complexes.
\end{hint}
\begin{solution}
Localization and graded structures refine Ext in common Noetherian settings. In this chapter the concrete consequence is that Ext measures the failure of Hom to be exact and classifies extension phenomena in low degree.
\end{solution}

\section*{Chapter summary}
The chapter now contributes a worked-problem layer to the progression from rings and modules through localization, Noetherian methods, integral dependence, and derived functors.
% BEGIN VOL05-EXPANSION V27-exercises-01
\section{Graded supplementary exercises}
These exercises extend the chapter with a balanced set of computations, proofs, hypothesis tests, applications, and challenges.

\subsection*{Standard computations}

\begin{exercise}[Cyclic Ext]\label{exr:v27-expand-01}
Compute $\operatorname{Ext}^1_{\mathbb Z}(\mathbb Z/6,\mathbb Z/15)$.
\end{exercise}
\begin{hint}
Use Hom on projective resolutions and extension classes.
\end{hint}
\begin{solution}
It is $\mathbb Z/3$.
\end{solution}

\begin{exercise}[Coprime Ext]\label{exr:v27-expand-02}
Compute $\operatorname{Ext}^1_{\mathbb Z}(\mathbb Z/4,\mathbb Z/9)$.
\end{exercise}
\begin{hint}
Use Hom on projective resolutions and extension classes.
\end{hint}
\begin{solution}
It is zero.
\end{solution}

\begin{exercise}[Projective]\label{exr:v27-expand-03}
Compute positive Ext from a free module.
\end{exercise}
\begin{hint}
Use Hom on projective resolutions and extension classes.
\end{hint}
\begin{solution}
It vanishes.
\end{solution}

\begin{exercise}[Degree zero]\label{exr:v27-expand-04}
What is $\operatorname{Ext}^0_R(M,N)$?
\end{exercise}
\begin{hint}
Use Hom on projective resolutions and extension classes.
\end{hint}
\begin{solution}
It is $\operatorname{Hom}_R(M,N)$.
\end{solution}

\begin{exercise}[Split class]\label{exr:v27-expand-05}
Which Ext class corresponds to a split short exact sequence?
\end{exercise}
\begin{hint}
Use Hom on projective resolutions and extension classes.
\end{hint}
\begin{solution}
The zero class.
\end{solution}

\subsection*{Proofs}

\begin{exercise}[Well-definedness of Ext proof]\label{exr:v27-expand-06}
Prove: $H^n(\operatorname{Hom}(P_\bullet,N))$ is independent of the chosen projective resolution of $M$.
\end{exercise}
\begin{hint}
Start from the theorem hypotheses and identify the decisive algebraic construction.
\end{hint}
\begin{solution}
Comparison maps between resolutions are unique up to chain homotopy. Applying Hom turns chain homotopies into cochain homotopies, which induce identical cohomology maps.
\end{solution}

\begin{exercise}[Projective vanishing proof]\label{exr:v27-expand-07}
Prove: If $P$ is projective, then $\operatorname{Ext}_A^n(P,N)=0$ for $n>0$.
\end{exercise}
\begin{hint}
Start from the theorem hypotheses and identify the decisive algebraic construction.
\end{hint}
\begin{solution}
Use the projective resolution concentrated in degree zero. The resulting Hom cochain complex has no positive cohomology.
\end{solution}

\begin{exercise}[Ext one classifies extensions proof]\label{exr:v27-expand-08}
Prove: Equivalence classes of short exact sequences $0\to N\to E\to M\to0$ form a group naturally isomorphic to $\operatorname{Ext}_A^1(M,N)$.
\end{exercise}
\begin{hint}
Start from the theorem hypotheses and identify the decisive algebraic construction.
\end{hint}
\begin{solution}
Pull an extension back along a projective presentation of $M$ and choose a splitting over the projective term. The resulting map from the first syzygy to $N$ defines a degree-one cohomology class; changing the splitting changes it by a coboundary, and the construction is reversible via pushout.
\end{solution}

\begin{exercise}[Structural synthesis]\label{exr:v27-expand-09}
Prove a corollary that genuinely uses both Well-definedness of Ext and Projective vanishing.
\end{exercise}
\begin{hint}
Apply the first theorem to obtain its structural description, then verify the second theorem's hypotheses.
\end{hint}
\begin{solution}
A correct proof explicitly invokes both results, verifies the common hypotheses, and derives a new consequence rather than merely restating either theorem.
\end{solution}

\subsection*{Counterexamples and hypothesis tests}

\begin{exercise}[Always split]\label{exr:v27-expand-10}
Test the claim: Every short exact sequence of modules splits.
\end{exercise}
\begin{hint}
Use the smallest concrete ring, module, or resolution that can decide the claim.
\end{hint}
\begin{solution}
False; nonzero Ext classes obstruct splitting.
\end{solution}

\begin{exercise}[Projective first variable]\label{exr:v27-expand-11}
Test the claim: A projective first argument can have nonzero positive Ext.
\end{exercise}
\begin{hint}
Use the smallest concrete ring, module, or resolution that can decide the claim.
\end{hint}
\begin{solution}
False.
\end{solution}

\begin{exercise}[Resolution dependence]\label{exr:v27-expand-12}
Test the claim: Ext depends on arbitrary choices in a projective resolution.
\end{exercise}
\begin{hint}
Use the smallest concrete ring, module, or resolution that can decide the claim.
\end{hint}
\begin{solution}
False up to canonical isomorphism.
\end{solution}

\subsection*{Applications and investigations}

\begin{exercise}[Extension obstruction]\label{exr:v27-expand-13}
What does $\operatorname{Ext}^1$ measure?
\end{exercise}
\begin{hint}
Translate the question into Hom on projective resolutions and extension classes.
\end{hint}
\begin{solution}
It measures equivalence classes of extensions and obstruction to splitting.
\end{solution}

\begin{exercise}[Deformation flavor]\label{exr:v27-expand-14}
Why does Ext occur in deformation theory?
\end{exercise}
\begin{hint}
Translate the question into Hom on projective resolutions and extension classes.
\end{hint}
\begin{solution}
Extension groups encode first-order gluing and obstruction data in many algebraic settings.
\end{solution}

\subsection*{Challenge problems}

\begin{exercise}[Local-global synthesis]\label{exr:v27-expand-15}
Connect 'Definition from projectives' with 'Duality preview' in one rigorous argument.
\end{exercise}
\begin{hint}
State the relevant ring map, module map, or complex before applying the structural theorem.
\end{hint}
\begin{solution}
The opening idea is: Choose a projective resolution $P_\bullet\to M$ and set $\operatorname{Ext}_A^n(M,N)=H^n(\operatorname{Hom}_A(P_\bullet,N))$. The later viewpoint is: Ext groups appear in duality, canonical modules, and deformation theory. A complete solution identifies the structural theorem that links these descriptions and checks its hypotheses.
\end{solution}

\begin{exercise}[Hypothesis audit]\label{exr:v27-expand-16}
Choose one theorem from the chapter, remove one essential hypothesis, and exhibit or explain a concrete failure.
\end{exercise}
\begin{hint}
Use one of the chapter's explicit computations or counterexamples as the test case.
\end{hint}
\begin{solution}
The solution must name the removed hypothesis, show exactly which proof step breaks, and give a concrete algebraic object where the claimed conclusion fails.
\end{solution}

% END VOL05-EXPANSION V27-exercises-01
